一个\emph{集合}是对象的汇集，这些对象称为该集合的\emph{元素}。
如果~$x$ 是~$A$ 的元素，就记作 $x \in A$。集合具有\emph{外延性}——它们完全由其元素确定。
我们可以通过\emph{显式列出}其元素，或者通过给出这些元素共有的性质来给定一个集合（\emph{抽象}）。外延性意味着：列出或给定元素的顺序与方式无关紧要。要证明 $A$ 和~$B$ 是同一个集合（$A = B$），就必须证明~$X$ 的每个元素都是~$Y$ 的元素，反之亦然。

重要的集合包括自然数集（$\Nat$）、整数集（$\Int$）、有理数集（$\Rat$）和实数集（$\Real$），也包括对象的\emph{字符串}集（$X^*$）和无穷\emph{序列}集（$X^\omega$）。
如果~$A$ 的每个元素也都是~$B$ 的元素，则 $A$ 是~$B$ 的\emph{子集}，记作 $A \subseteq B$。
一个集合~$B$ 的所有子集构成的汇集本身也是一个集合，称为~$B$ 的\emph{幂集} $\Pow{B}$。
我们可以构造集合的\emph{并集} $A \cup B$ 和\emph{交集} $A \cap B$。
一个\emph{有序对} $\tuple{x, y}$ 由两个对象~$x$ 和~$y$ 依此特定顺序组成。有序对 $\tuple{x, y}$ 和 $\tuple{y, x}$ 是不同的有序对（除非 $x = y$）。
所有满足 $x \in A$ 且 $y \in B$ 的有序对 $\tuple{x, y}$ 构成的集合，称为 $A$ 与~$B$ 的\emph{笛卡尔积}~$A \times B$。
我们用 $A^2$ 表示 $A \times A$；例如 $\Nat^2$ 就是自然数有序对的集合。